%% kapa: exam=1 round=34 year=2023 period=2 form=A subject=appraisal_law subject_ko=감정평가관계법규 seq=33 number=33
%% source: Q-Net 감정평가사 기출문제 (cst00309) artlSeq=5211789 file=raw/1008/5211789/2023년도 제34회 감정평가사 1차시험 2교시.pdf
%% conversion: kordoc 4.24.51 -> tools/gen1.py (auto); review=proofread-sonnet; flags=none
\begin{question}{33}
부동산등기법령상 권리에 관한 등기에 관한 설명으로 옳지 않은 것은?
\begin{choices}
  \item 국가·지방자치단체·국제기관 및 외국정부의 부동산등기용등록번호는 법무부장관이 지정·고시한다.
  \item 등기관이 환매특약의 등기를 할 때에는 매수인이 지급한 대금을 기록하여야 한다.
  \item 등기원인에 권리의 소멸에 관한 약정이 있을 경우 신청인은 그 약정에 관한 등기를 신청할 수 있다.
  \item 등기의 말소를 신청하는 경우에 그 말소에 대하여 등기상 이해관계 있는 제3자가 있을 때에는 제3자의 승낙이 있어야 한다.
  \item 등기관이 토지에 대하여 소유권경정등기를 하였을 때에는 지체 없이 그 사실을 지적소관청에 알려야 한다.
\end{choices}
\end{question}
\answer{1}
